\section{Roth: the effective boundary and a complete quadratic theorem}
\label{sec:roth}

\subsection{Three different meanings of finiteness}

For a real algebraic irrational $\alpha$ and rational $\delta=a/b>0$, where
$a,b\geq1$ are integers, Roth's theorem says that the reduced fractions
$p/q$ satisfying
\begin{equation}\label{eq:roth-inequality}
  \left|\alpha-\frac pq\right|<q^{-2-a/b},\qquad q\geq1,
\end{equation}
form a finite set.  A constructive formulation must specify the data meant by
``finite.''  The following are different assertions.
\begin{enumerate}
\item A \emph{cutoff}: produce an integer $Q$ such that every solution has
$q<Q$.  This permits a finite exhaustive search.
\item A \emph{cardinality bound}: produce $N$ and prove that no list of $N+1$
distinct reduced fractions consists entirely of solutions.
\item A \emph{finite obstruction}: for every supplied finite list satisfying
specified growth conditions, locate a member that is not a solution.
\end{enumerate}
A cutoff is not supplied by an upper bound on cardinality.  Indeed, let
$P_e(n)$ say that program $e$ halts for the first time at step $n$.  This is a
primitive recursive predicate, and for every $e$ there is at most one $n$
with $P_e(n)$.  A computable function $B(e)$ bounding all such $n$ would
solve the halting problem by simulating $e$ for $B(e)$ steps.  Thus no general
algorithm converts cardinality bounds into cutoffs.

Roth's theorem is known to have effective cardinality estimates; Luckhardt's
analysis of the Roth and Esnault--Viehweg proofs is a particularly relevant
example \cite{Luckhardt1989,KohlenbachHerbrand}.  The general cutoff problem
is a different, unresolved effectivity problem.  The introduction of
Adiceam and Shirandami's 2026 paper explicitly describes the corresponding
height problem for the Subspace Theorem, including Roth's case, as open
\cite{AdiceamShirandami2026}.  Their probabilistic results do not provide a
cutoff for each specified algebraic number.

Consequently, a uniform witness-producing proof of
\begin{equation}\label{eq:effective-roth-target}
 \forall A,a,b\ \exists Q\ \forall p,q\quad
 \bigl(\operatorname{Valid}(A)\wedge a,b,q\geq1
       \wedge\operatorname{Good}_{A,a,b}(p,q)\bigr)\Longrightarrow q<Q
\end{equation}
would settle the effective Roth problem, if $A$ ranges over finite algebraic
number codes and the premise is arranged in the usual decidable-input form.
Here and below one can equivalently move the input validity premises outside
the existential quantifier.  A constructive proof in the intended arithmetic
would supply a computable bound; under the explicit primitive recursive
interpretation required in this report it must supply a primitive recursive
one.  This observation is a barrier to claiming the requested full theorem,
\emph{not} a proof that no such proof exists.  Nor does it rule out an iPRA
proof of a suitable quantitative finite-obstruction version.  Such a proof
would require a complete arithmetic treatment of Roth's auxiliary-polynomial
and index estimates.  Those estimates are not proved in this report and are
not smuggled into an axiom named ``effective Roth.''

We next prove an unconditional theorem that does have the required explicit
arithmetic content: the full Roth cutoff for quadratic algebraic irrationals,
and the Liouville range for every degree.  We give both the familiar
human-readable argument and its entirely rational, finitely checkable form.

\subsection{Finite input data and the separation constant}

Let
\[
 f(X)=\sum_{i=0}^{d}c_iX^i\in\mathbb Z[X],\qquad d\geq2,\quad c_d\ne0.
\]
Assume $f$ has no rational root.  This assumption admits a finite certificate.
Require $c_0\ne0$ and check
\begin{equation}\label{eq:no-root-certificate}
 \sum_{i=0}^{d}c_i u^i v^{d-i}\ne0
 \quad\text{for every } |u|\leq |c_0|,\quad 1\leq v\leq |c_d|.
\end{equation}
There is no need to restrict the finite check to coprime pairs.  The
rational-root theorem proves its sufficiency: if a reduced $p/q$ is a root,
reduction of $q^df(p/q)=0$ modulo $p$ and modulo $q$, followed by Euclid's
lemma, gives $p\mid c_0$ and $q\mid c_d$.  The case $p=0$ is excluded by
$c_0\ne0$.  Therefore $|p|\leq|c_0|$ and $q\leq|c_d|$, contradicting the
check.  The same conclusion then holds for unreduced fractions by reducing
them.  Euclid's lemma here follows from the integer extended Euclidean
algorithm, or directly from its finite Bezout certificate.

Take rational numbers $u_0<v_0$ with $f(u_0)f(v_0)<0$, and an integer $R\geq1$
with $|u_0|,|v_0|\leq R$.  These are all decidable rational conditions.  Put
\begin{equation}\label{eq:roth-K}
 K=\max\left(1,\sum_{i=1}^{d} i|c_i|(R+1)^{i-1}\right).
\end{equation}
For rational or real $x,y$ with $|x|,|y|\leq R+1$, the identity
\[
 x^i-y^i=(x-y)\sum_{j=0}^{i-1}x^{i-1-j}y^j
\]
gives
\begin{equation}\label{eq:polynomial-lipschitz}
 |f(x)-f(y)|\leq K|x-y|.
\end{equation}
This proof uses only finite sums, the triangle inequality, and multiplication;
no derivative, mean-value theorem, compactness principle, or maximum is
needed.

The input defines a primitive recursive rational name $r_k$.  Starting from
$[u_0,v_0]$, bisect $2R+k$ times, retaining the half whose endpoint values
have opposite signs, and let $r_k$ be the left endpoint.  Every midpoint is
rational and cannot be a root by \eqref{eq:no-root-certificate}, so the
branch is decided by an integer sign computation.  Write $s_k$ for the right
endpoint after the same number of steps.  Since
$v_0-u_0\leq2R\leq2^{2R}$, the construction gives
\begin{equation}\label{eq:root-name-properties}
 \begin{split}
 &|r_k|,|s_k|\leq R,\qquad 0<s_k-r_k\leq2^{-k},\\
 &|r_k-r_l|\leq2^{-k}+2^{-l},\qquad |f(r_k)|\leq K2^{-k}.
 \end{split}
\end{equation}
For the last inequality, opposite signs imply
$|f(r_k)|\leq|f(r_k)-f(s_k)|$, and then
\eqref{eq:polynomial-lipschitz} applies.  For the Cauchy inequality, when
$k\leq l$ the later interval lies inside the earlier one, so the difference
between left endpoints is at most $2^{-k}$; the other order is symmetric.
These are finite induction arguments about a bounded bisection computation.
The symbol $\alpha$ below denotes the real described by this name; the
arithmetic proof itself needs only \eqref{eq:root-name-properties}.

\subsection{Liouville separation without limits}

\begin{theorem}[Explicit algebraic separation]\label{thm:liouville-explicit}
For every integer $p$, positive integer $q$, and $k\geq0$,
\begin{equation}\label{eq:finite-separation}
 \left|\frac pq-r_k\right|+2^{-k}\ \geq\ \frac1{Kq^d}.
\end{equation}
In the conventional real notation this implies
\begin{equation}\label{eq:liouville-real}
 \left|\frac pq-\alpha\right|\geq\frac1{Kq^d}.
\end{equation}
\end{theorem}

\begin{proof}
Let $t=p/q$.  The number
\[
 N_f(p,q)=q^df(p/q)=\sum_{i=0}^{d}c_i p^i q^{d-i}
\]
is a nonzero integer, and therefore
\begin{equation}\label{eq:integer-gap}
 |f(t)|\geq q^{-d}.
\end{equation}
The rational comparison $|p|>(R+1)q$ is decidable.  If it holds, then
$|t-r_k|>1$, while $1/(Kq^d)\leq1$, so the assertion follows.
In the remaining case $|t|\leq R+1$.  For any $l$, equations
\eqref{eq:polynomial-lipschitz}, \eqref{eq:root-name-properties}, and
\eqref{eq:integer-gap} give
\[
 q^{-d}\leq|f(t)|
 \leq |f(t)-f(r_l)|+|f(r_l)|
 \leq K|t-r_l|+K2^{-l}.
\]
Taking $l=k$ and dividing by the positive integer $K$ proves
\eqref{eq:finite-separation} immediately.  Every comparison in this
argument is rational and decidable.
The real formulation is its usual Cauchy-name interpretation.
\end{proof}

For a reader working directly with real algebraic numbers, there is a shorter
presentation of the same argument.  If $|p|\leq(R+1)q$, then $f(\alpha)=0$
and \eqref{eq:polynomial-lipschitz} imply
$q^{-d}\leq|f(p/q)|\leq K|p/q-\alpha|$; the complementary rational branch
again gives an error exceeding $1$.  The expanded proof above explains why
this familiar argument does not conceal an infinitary principle.

\subsection{An explicit Roth cutoff for every quadratic input}

Fractional powers need not appear in the arithmetic language.  Define a
\emph{positive certificate of Roth approximation} to be a triple $(p,q,k)$
with $p\in\mathbb Z$, $q\geq1$, and
\begin{equation}\label{eq:good-certificate}
 q^{2b+a}\left(\left|\frac pq-r_k\right|+2^{-k}\right)^b<1.
\end{equation}
This is a decidable rational inequality, and clearing denominators makes it
an integer inequality.  It asserts a strict upper bound with an explicit
error allowance.  Conversely, a strict real inequality
\eqref{eq:roth-inequality} supplies such a certificate once a positive
rational gap and a sufficient Cauchy precision are given.  No trichotomy
principle for arbitrary real numbers is being assumed.

\begin{theorem}[Primitive recursive cutoff in the Liouville range]
\label{thm:quadratic-roth}
Let $a,b\geq1$ and suppose the integer
\[
 c=a-(d-2)b
\]
is positive.  Every certificate \eqref{eq:good-certificate} satisfies
\begin{equation}\label{eq:cutoff-explicit}
 q^c<K^b,\qquad\text{and hence}\qquad q<K^b.
\end{equation}
In particular, when $d=2$, this proves the full Roth finiteness conclusion
with the explicit cutoff $Q=K^b$ for every positive rational $\delta=a/b$.
\end{theorem}

\begin{proof}
Write $E=|p/q-r_k|+2^{-k}$.  The preceding theorem gives
$E\geq1/(Kq^d)$.  Thus \eqref{eq:good-certificate} implies
\[
 1>q^{2b+a}E^b
   \geq\frac{q^{2b+a}}{K^bq^{db}}
   =\frac{q^c}{K^b}.
\]
As $c\geq1$ and $q\geq1$, finite induction gives $q\leq q^c$.
This proves both inequalities.  For $d=2$, $c=a\geq1$ automatically.
\end{proof}

All approximants are therefore found among the finite pairs
\[
 1\leq q<K^b,\qquad |p|\leq (R+1)q.
\]
The numerator bound follows because a good approximation has error less
than $1$.  Exact membership can be decided from the algebraic input by
sign determination, but even without implementing sign determination the
finite rectangle is already the promised cutoff certificate.  Coprimality
can be checked by the Euclidean algorithm; the proof of the bound did not
require it.

For degree $d\geq3$, the hypothesis is precisely
$\delta>d-2$.  It is essential to identify this as the Liouville range,
rather than to call it general Roth.  The difficult range
$0<\delta\leq d-2$ is not settled by this argument.  For example, the input
$f(X)=X^3-2$ with $0<\delta\leq1$ is outside the theorem's scope.

\paragraph{Concrete example.}
For $f(X)=X^2-2$, take $[u_0,v_0]=[1,2]$ and $R=2$.
The finite rational-root check succeeds, and \eqref{eq:roth-K} gives $K=6$.
Thus every fraction satisfying
$|\sqrt2-p/q|<q^{-2-a/b}$ has $q<6^b$.
For $\delta=1/2$, the explicit cutoff is $q<36$.
These deliberately generous constants make the complete arithmetic proof
transparent; sharper quadratic constants are unnecessary for the theorem.

\subsection{What the elementary gap principle does and does not give}

Suppose two distinct reduced fractions $p/q$ and $r/s$, with $1\leq q\leq s$,
both satisfy \eqref{eq:roth-inequality}.  Since $ps-rq$ is a nonzero integer,
\[
 \frac1{qs}\leq\left|\frac pq-\frac rs\right|
 <q^{-2-\delta}+s^{-2-\delta}\leq2q^{-2-\delta},
\]
and hence
\begin{equation}\label{eq:gap-principle}
 s>\tfrac12q^{1+\delta}.
\end{equation}
For $\delta=a/b>0$ and $q\geq4^b$, this implies $s>2q$.
The identical argument can be conducted with rational positive certificates
and integer powers.  The gap principle forbids clustering and makes the
first successive denominators grow rapidly.  It does not exclude an
arbitrarily sparse infinite sequence; a global finite bound requires the
additional Roth estimates.

There is nevertheless useful constructive content in a certified
cardinality bound.  Suppose a decidable predicate $G(n)$ has at most $N$
solutions, and let $F$ be a supplied primitive recursive function with
$F(t)\geq t$.  Define
\[
 B_0=1,\qquad B_{i+1}=F(B_i)+1\quad(0\leq i\leq N).
\]
The $N+1$ disjoint finite intervals $[B_i,B_{i+1})$ cannot each contain a
solution.  Bounded searches find an empty one.  Thus there is an explicit
primitive recursive bound, depending on the supplied $F$, on a location of
a finite window with no solutions.  This is a valid route from a
quantitative Roth theorem to a metastable formulation, provided the
cardinality theorem and its arithmetic proof are actually supplied.
It still supplies no final cutoff valid for all later denominators.

\subsection{Lean and object-theory proof obligations}

The proved fragment can be organized into the following modules.
\begin{enumerate}
\item \texttt{RationalArithmetic}: signed integer and positive-denominator
codes; addition, multiplication, absolute value, powers, and decidable order.
Expose the cleared-denominator correctness lemmas.
\item \texttt{PolynomialDifference}: prove the finite factorization of
$x^i-y^i$ by induction on $i$, then
\eqref{eq:polynomial-lipschitz} by induction over the coefficient list.
\item \texttt{NoRationalRoot}: verify \eqref{eq:no-root-certificate}; prove
that $N_f(p,q)\ne0$ for arbitrary $p,q$ by rational reduction and the
rational-root argument.  The Euclidean algorithm needs a declared bounded
implementation or an explicit primitive recursive majorant.
\item \texttt{RootBisection}: define the bounded bisection loop on rational
endpoints and prove nesting, sign preservation, width, and residual bounds.
\item \texttt{LiouvilleFinite}: prove \eqref{eq:finite-separation} directly from the integer gap,
the difference bound, and the bisection residual bound at the same precision.
\item \texttt{QuadraticRoth}: define \eqref{eq:good-certificate} and prove
\eqref{eq:cutoff-explicit} by integer power arithmetic.
\end{enumerate}

The object-theoretic statements are schemes for fixed primitive recursive
operations, or use one fixed primitive recursive evaluator for the stated
finite polynomial and bisection data.  They do not posit a primitive
recursive evaluator for arbitrary primitive recursive programs.  Bounded
quantifiers in the input checks are implemented by Boolean folds; induction
in the argument is on explicit finite loops and can be presented as
quantifier-free induction on their invariants after coding finite records.
One must still translate the resulting derivations into the repository's
chosen iPRA calculus.  A Lean theorem over \texttt{Real} or a successful
\texttt{ring}/\texttt{norm\_num} invocation, by itself, would not accomplish
that translation.

For the full Roth target, a formalization plan must list the auxiliary
polynomial construction, coefficient estimates, index bounds, quantitative
selection of approximants, and extraction of the selected constructive
conclusion as separate proof obligations.  Calling any one of these an
unproved axiom would produce a conditional development.  A complete Lean
proof of the cutoff version would first require solving the mathematical
effectivity problem; the present report claims neither that result nor its
formalization.
